%!TEX root = ../Thesis.tex
\chapter{District Heating System}

\section{Derivation of $\dot{T}_i$ Equation}
\label{sec:district_heating_temperature_derivation}

The district heating pipes are assumed to transport incompressible water with
constant specific heat capacity. The flow velocity is assumed to be constant
throughout each pipe, and only temperature variations in the axial direction
are considered.

For a pipe with axial coordinate \(z\), the energy balance is written as

\begin{equation}
    \partial_t u(t,z)
    =
    -
    \partial_z \mathcal{N}_h(t,z)
    +
    q(t,z),
    \qquad
    z \in [z_0,z_f]
\end{equation}

where \(u(t,z)\) denotes the internal energy per unit volume,
\(\mathcal{N}_h(t,z)\) denotes the enthalpy flux, and \(q(t,z)\) represents
the heat loss per unit volume.

The enthalpy flux is expressed as

\begin{equation}
    \mathcal{N}_h(t,z)
    =
    H
    \left(
        T(t,z),p,\mathcal{N}(t)
    \right)
\end{equation}

where \(T(t,z)\) is the water temperature, \(p\) is the pressure, and
\(\mathcal{N}(t)\) is the molar flux. The molar flux is given by

\begin{equation}
    \mathcal{N}(t)
    =
    \rho v(t)
\end{equation}

where \(\rho\) denotes the constant density and \(v(t)\) denotes the flow
velocity.

The heat loss to the surroundings is modeled as

\begin{equation}
    q(t,z)
    =
    -
    k
    \left(
        T(t,z)-T_a(t)
    \right)
\end{equation}

where \(k\) is the conductance and \(T_a(t)\) is the ambient temperature.

Applying the chain rule to the enthalpy flux gives

\begin{equation}
    \partial_T
    H
    \left(
        T(t,z),p,\mathcal{N}(t)
    \right)
    \partial_z T(t,z)
    =
    \partial_z \mathcal{N}_h(t,z)
\end{equation}

Using the specific heat capacity \(c_p\), this becomes

\begin{equation}
    \mathcal{N}(t)c_p
    \partial_z T(t,z)
    =
    \partial_z \mathcal{N}_h(t,z)
\end{equation}

The pipe is then segmented into \(M\) finite volumes. For the \(i\)-th
segment,

\begin{equation}
    \Omega_i
    =
    [z_i,z_{i+1}]
\end{equation}

and the energy balance is integrated over the segment:

\begin{equation}
    \partial_t
    \int_{\Omega_i}
    u(t,z)\,dz
    =
    -
    \int_{\Omega_i}
    \partial_z \mathcal{N}_h(t,z)\,dz
    +
    \int_{\Omega_i}
    q(t,z)\,dz
\end{equation}

Multiplying by the cross-sectional area \(A\), the energy balance of the
segment can be written as

\begin{equation}
    \dot{U}_i(t)
    =
    -
    A\mathcal{N}(t)c_p
    \int_{z_i}^{z_{i+1}}
    \partial_z T(t,z)\,dz
    +
    Q_i(t)
\end{equation}

where \(U_i(t)\) denotes the internal energy of the segment and \(Q_i(t)\)
denotes the corresponding heat loss.

Evaluating the integral gives

\begin{equation}
    \dot{U}_i(t)
    =
    -
    A\mathcal{N}(t)c_p
    \left(
        T(t,z_{i+1})
        -
        T(t,z_i)
    \right)
    +
    Q_i(t)
\end{equation}

Using the average temperature of each finite volume, this is approximated as

\begin{equation}
    \dot{U}_i(t)
    =
    -
    A\mathcal{N}(t)c_p
    \left(
        \bar{T}_{i+1}(t)
        -
        \bar{T}_i(t)
    \right)
    +
    Q_i(t)
\end{equation}

The internal energy of the \(i\)-th segment is written as

\begin{equation}
    U_i(t)
    =
    n c_p \bar{T}_i(t)
\end{equation}

where

\begin{equation}
    n
    =
    A\Delta z \rho,
    \qquad
    \Delta z
    =
    z_{i+1}-z_i
\end{equation}

The heat loss of the segment is approximated by

\begin{equation}
    Q_i(t)
    \approx
    -
    A\Delta z k
    \left(
        \bar{T}_i(t)-T_a(t)
    \right)
\end{equation}

Substituting the internal-energy and heat-loss expressions into the energy
balance gives

\begin{align}
    c_p A\Delta z \rho \,
    \dot{\bar{T}}_i(t)
    &=
    -
    A\mathcal{N}(t)c_p
    \left(
        \bar{T}_{i+1}(t)
        -
        \bar{T}_i(t)
    \right)
    \nonumber\\
    &\quad
    -
    A\Delta z k
    \left(
        \bar{T}_i(t)-T_a(t)
    \right)
\end{align}

Since

\begin{equation}
    \mathcal{N}(t)
    =
    \rho v(t)
\end{equation}

the equation can be divided by \(A\Delta z\rho c_p\), resulting in

\begin{equation}
    \boxed{
    \dot{\bar{T}}_i(t)
    =
    -
    \frac{v(t)}{\Delta z}
    \left(
        \bar{T}_{i+1}(t)
        -
        \bar{T}_i(t)
    \right)
    -
    \frac{k}{\rho c_p}
    \left(
        \bar{T}_i(t)
        -
        T_a(t)
    \right)
    }
    \label{eq:district_heating_temperature}
\end{equation}

The resulting pipe model can subsequently be expressed in continuous
state-space form as

\begin{align}
    \dot{x}(t)
    &=
    A_c x(t)
    +
    B_c u(t)
    +
    E_c d(t)\\
    y(t)
    &=
    C_c x(t)
    +
    e(t_k)
\end{align}

where the states \(x(t)\) correspond to the average temperatures
\(\bar{T}_i(t)\), \(u(t)\) represents the supply temperature \(T_s\),
\(d(t)\) represents the ambient temperature \(T_a\), and \(y(t)\) contains
the temperatures measured at the selected pipe nodes.


\section{Change of Variable}
\label{sec:district_heating_change_variable}

From the temperature equation derived in the previous section,

\begin{equation}
    \dot{T}_i
    =
    -
    \frac{v}{\Delta z}
    \left(
        T_i-T_{i-1}
    \right)
    -
    \frac{k}{\rho c_p}
    \left(
        T_i-T_a
    \right)
\end{equation}

To simplify the notation, define

\begin{equation}
    \alpha(v)
    =
    \frac{v}{\Delta z}
\end{equation}

and

\begin{equation}
    \beta
    =
    \frac{k}{\rho c_p}
\end{equation}

The temperature dynamics can then be written as

\begin{equation}
    \dot{T}_i
    =
    -
    \alpha(v)
    \left(
        T_i-T_{i-1}
    \right)
    -
    \beta
    \left(
        T_i-T_a
    \right)
    \label{eq:temperature_alpha_beta}
\end{equation}

A change of variables is introduced by defining the state relative to the
ambient temperature as

\begin{equation}
    x_i
    =
    T_i-T_a
    \label{eq:change_state}
\end{equation}

Similarly, the input is defined as the supply temperature relative to the
ambient temperature,

\begin{equation}
    u
    =
    T_s-T_a
    \label{eq:change_input}
\end{equation}

Therefore,

\begin{equation}
    T_i
    =
    x_i+T_a
\end{equation}

and

\begin{equation}
    T_s
    =
    u+T_a
\end{equation}

Differentiating \cref{eq:change_state} gives

\begin{equation}
    \dot{x}_i
    =
    \dot{T}_i-\dot{T}_a
\end{equation}

Substituting

\begin{equation}
    T_i=x_i+T_a,
    \qquad
    T_{i-1}=x_{i-1}+T_a
\end{equation}

into \cref{eq:temperature_alpha_beta} gives

\begin{equation}
    \dot{x}_i+\dot{T}_a
    =
    -
    \alpha
    \left[
        (x_i+T_a)
        -
        (x_{i-1}+T_a)
    \right]
    -
    \beta x_i
\end{equation}

After simplifying,

\begin{equation}
    \boxed{
    \dot{x}_i
    =
    -
    (\alpha+\beta)x_i
    +
    \alpha x_{i-1}
    -
    \dot{T}_a
    }
    \label{eq:district_heating_xi}
\end{equation}

For the first pipe segment, the upstream temperature is the supply
temperature \(T_s\). Since

\begin{equation}
    u=T_s-T_a
\end{equation}

the first state equation becomes

\begin{equation}
    \boxed{
    \dot{x}_0
    =
    -
    (\alpha+\beta)x_0
    +
    \alpha u
    -
    \dot{T}_a
    }
    \label{eq:district_heating_x0}
\end{equation}

Collecting all pipe-segment temperatures into the state vector gives the
continuous-time system

\begin{equation}
    \boxed{
    \dot{x}
    =
    A_c(v)x
    +
    B_c(v)u
    +
    E_c\dot{T}_a
    }
    \label{eq:district_heating_continuous}
\end{equation}

The system matrix is

\begin{equation}
A_c(v)
=
\begin{bmatrix}
-(\alpha+\beta) & 0 & 0 & \cdots & 0 \\
\alpha & -(\alpha+\beta) & 0 & \cdots & 0 \\
0 & \alpha & -(\alpha+\beta) & \ddots & \vdots \\
\vdots & \ddots & \ddots & \ddots & 0 \\
0 & \cdots & 0 & \alpha & -(\alpha+\beta)
\end{bmatrix}
\end{equation}

where

\begin{equation}
    \alpha=\alpha(v)=\frac{v}{\Delta z}
\end{equation}

The input matrix is

\begin{equation}
B_c(v)
=
\begin{bmatrix}
\alpha \\
0 \\
\vdots \\
0
\end{bmatrix}
\end{equation}

and the disturbance matrix associated with the change in ambient temperature
is

\begin{equation}
E_c
=
\begin{bmatrix}
-1 \\
-1 \\
\vdots \\
-1
\end{bmatrix}
\end{equation}

\section{Controller and Hidden Plant Approach}
\label{sec:controller_hidden_plant}

After introducing the change of variable

\begin{equation}
    x_i = T_i-T_a
\end{equation}

the continuous-time district heating model was obtained as

\begin{equation}
    \dot{x}
    =
    A_c(v)x
    +
    B_c(v)u
    +
    E_c\dot{T}_a
    \label{eq:district_full_continuous}
\end{equation}

The change of variable therefore introduces the derivative of the ambient
temperature into the state dynamics. In the implementation, the model used
by the controller and the model used to simulate the physical plant are
treated separately.


\subsection{Controller Approach}
\label{subsec:district_controller}

On the controller side, the ambient temperature is assumed to remain constant
within each sampling interval. Thus,

\begin{equation}
    T_a(t)
    =
    T_{a,k},
    \qquad
    t\in[t_k,t_{k+1})
\end{equation}

and consequently

\begin{equation}
    \dot{T}_a(t)
    =
    0,
    \qquad
    t\in(t_k,t_{k+1})
\end{equation}

Therefore,

\begin{equation}
    E_c\dot{T}_a=0
\end{equation}

and the continuous-time model used by the controller reduces to

\begin{equation}
    \boxed{
    \dot{x}
    =
    A_c(v)x
    +
    B_c(v)u
    }
    \label{eq:district_controller_continuous}
\end{equation}

This model is discretized using a Zero-Order Hold assumption. During each
sampling interval, the matrices and the input are kept constant. The
resulting state at the end of the sampling interval is

\begin{equation}
    x(t_{k+1}^{-})
    =
    e^{A_{c,k}\Delta t}x_k
    +
    \int_{0}^{\Delta t}
    e^{A_{c,k}(\Delta t-\tau)}
    B_{c,k}u_k
    \,d\tau
    \label{eq:district_zoh}
\end{equation}

The discrete-time matrices are obtained from

\begin{equation}
    \exp
    \left(
    \begin{bmatrix}
        A_c & B_c \\
        0 & 0
    \end{bmatrix}
    \Delta t
    \right)
    =
    \begin{bmatrix}
        A_k & B_k \\
        0 & I
    \end{bmatrix}
    \label{eq:district_zoh_matrix}
\end{equation}

Thus, the controller uses the discrete-time model generated from the
continuous dynamics under the assumption that the ambient temperature does
not change inside the sampling interval.


\subsection{Hidden Plant Approach}
\label{subsec:district_hidden_plant}

The hidden plant represents the system used to generate the simulated
measurements. In contrast to the controller model, the hidden plant is
affected by the actual ambient-temperature values.

Because the transformed state is defined relative to the ambient temperature,
a change in \(T_a\) between two sampling instants changes the state
coordinates. After the continuous evolution over one sampling interval,

\begin{equation}
    x_{k+1}
    =
    x_{k+1}^{-}
    +
    \left(
        T_{a,k}-T_{a,k+1}
    \right)\mathbf{1}
    \label{eq:district_ambient_shift}
\end{equation}

The ambient-temperature contribution is therefore defined as

\begin{equation}
    d_k
    =
    \left(
        T_{a,k}-T_{a,k+1}
    \right)\mathbf{1}
    \label{eq:district_dk}
\end{equation}

The hidden plant simulation consequently uses

\begin{equation}
    \boxed{
    x_{k+1}
    =
    A_kx_k
    +
    B_ku_k
    +
    d_k
    +
    w_k
    }
    \label{eq:district_hidden_state}
\end{equation}

where

\begin{equation}
    w_k
    \sim
    \mathcal{N}(0,R_w)
\end{equation}

The output represents the outlet temperature relative to the current ambient
temperature,

\begin{equation}
    y_k
    =
    T_{\mathrm{out},k}
    -
    T_{a,k}
\end{equation}

The measured output is modeled as

\begin{equation}
    \boxed{
    y_k
    =
    Cx_k+\nu_k
    }
    \label{eq:district_hidden_output}
\end{equation}

with

\begin{equation}
    \nu_k
    \sim
    \mathcal{N}(0,R_v)
\end{equation}

The hidden plant is used only to simulate the underlying system behavior.
Its internal state is not assumed to be directly available to the
controller. Instead, the controller receives the noisy output measurements,
and the Kalman filter is applied to these measurements in order to estimate
the internal state.

Therefore, the controller performs its prediction using the model in which
the ambient temperature is assumed constant over each sampling interval,
whereas the hidden plant evolves according to the actual ambient-temperature
sequence through the additional term \(d_k\). This separation allows the
effect of changing ambient conditions to be included in the simulation while
maintaining the controller formulation based on the ZOH-discretized model.

\subsection{Change on Kalman Filter}
\label{subsec:district_kalman_change}

The change of variable

\begin{equation}
    x_i=T_i-T_a
\end{equation}

also introduces a modification to the Kalman-filter procedure. Since the
state is defined relative to the ambient temperature, a change in the ambient
temperature between two consecutive sampling instants changes the state
coordinates.

As derived previously,

\begin{equation}
    x_{k+1}
    =
    x_{k+1}^{-}
    +
    \left(
        T_{a,k}-T_{a,k+1}
    \right)\mathbf{1}
\end{equation}

Defining

\begin{equation}
    d_k
    =
    \left(
        T_{a,k}-T_{a,k+1}
    \right)\mathbf{1}
\end{equation}

the same correction must also be applied to the predicted state estimate of
the Kalman filter.

The measurement-update equations remain unchanged. Starting from the prior
estimate \(\hat{x}_{k|k-1}\), the measurement \(y_k\) is used to obtain

\begin{equation}
    \hat{x}_{k|k}
    =
    \hat{x}_{k|k-1}
    +
    P_{k|k-1}C^T
    \left(
        CP_{k|k-1}C^T+R_v
    \right)^{-1}
    \left(
        y_k-C\hat{x}_{k|k-1}
    \right)
\end{equation}

and

\begin{equation}
    P_{k|k}
    =
    P_{k|k-1}
    -
    P_{k|k-1}C^T
    \left(
        CP_{k|k-1}C^T+R_v
    \right)^{-1}
    CP_{k|k-1}
\end{equation}

The time update is then first performed using the controller model,

\begin{equation}
    \hat{x}_{k+1|k}^{-}
    =
    A_k\hat{x}_{k|k}
    +
    B_ku_k
\end{equation}

At this stage, the predicted state is based on the assumption used by the
controller that the ambient temperature remains constant over the sampling
interval.

When the next ambient-temperature value \(T_{a,k+1}\) becomes available, the
change in ambient temperature is calculated through

\begin{equation}
    d_k
    =
    \left(
        T_{a,k}-T_{a,k+1}
    \right)\mathbf{1}
\end{equation}

The predicted state estimate is then corrected as

\begin{equation}
    \boxed{
    \hat{x}_{k+1|k}
    =
    \hat{x}_{k+1|k}^{-}
    +
    d_k
    }
\end{equation}

Therefore, the complete prediction step can be written as

\begin{equation}
    \boxed{
    \hat{x}_{k+1|k}
    =
    A_k\hat{x}_{k|k}
    +
    B_ku_k
    +
    d_k
    }
\end{equation}

After this correction, the time index is advanced,

\begin{equation}
    k \leftarrow k+1
\end{equation}

and the corrected prediction becomes the prior estimate used in the next
measurement update.

Thus, the main modification to the Kalman-filter structure is the additional
state correction introduced when the new ambient-temperature value becomes
available. This allows the state estimate to remain consistent with the
definition \(x_i=T_i-T_a\) while the actual ambient temperature changes over
time.

\section{Updates on Implementation}
\label{sec:district_heating_implementation_updates}

Two modifications were introduced in the implementation of the optimization
problem. First, the changes in the control input \(U\) and flow velocity
\(V\) were included in the cost function. Second, the numerical optimization
method was changed from SLSQP to IPOPT due to the increase in computational
time.

\subsection{Penalizing Changes in \(U\) and \(V\)}

Previously, the objective function penalized the tracking error together with
the magnitudes of the control input \(U\) and the flow velocity \(V\).
However, minimizing only these quantities does not explicitly discourage
large changes between consecutive control actions.

Therefore, additional terms were introduced to penalize the changes in
\(U\) and \(V\). For consecutive prediction steps, these changes are defined
as

\begin{equation}
    \Delta u_k
    =
    u_k-u_{k-1}
\end{equation}

and

\begin{equation}
    \Delta v_k
    =
    v_k-v_{k-1}
\end{equation}

The corresponding penalties are

\begin{equation}
    \Delta u_k^T R_{\Delta U}\Delta u_k
\end{equation}

and

\begin{equation}
    \Delta v_k^T R_{\Delta V}\Delta v_k
\end{equation}

The updated objective function therefore becomes

\begin{align}
    \phi(U,V)
    =
    \frac{1}{2}
    \sum_{k=0}^{N-1}
    \Big[
    &
    (\mu_{y,k}-\bar{y}_k)^T
    Q
    (\mu_{y,k}-\bar{y}_k)
    +
    u_k^T R_U u_k
    +
    v_k^T R_V v_k
    \nonumber\\
    &
    +
    \operatorname{tr}(Q P_{y,k})
    +
    \Delta u_k^T
    R_{\Delta U}
    \Delta u_k
    +
    \Delta v_k^T
    R_{\Delta V}
    \Delta v_k
    \Big]
    \label{eq:district_updated_cost}
\end{align}

Thus, \(R\) and \(S\) penalize the magnitude of the control variables,
whereas \(Q_{\mathrm{change}}\) and \(S_{\mathrm{change}}\) penalize how
rapidly these variables change from one time step to the next.

For \(k>0\),

\begin{align}
    \Delta u_k &= u_k-u_{k-1},\\
    \Delta v_k &= v_k-v_{k-1}
\end{align}

In the receding-horizon implementation, the first elements of the newly
optimized sequences are compared with the values that were actually applied
at the previous sampling instant. Therefore,

\begin{align}
    \Delta u_0
    &=
    u_0-u_{\mathrm{prev}},\\
    \Delta v_0
    &=
    v_0-v_{\mathrm{prev}}
\end{align}

When no previous control value exists, such as at the beginning of the
simulation, this first change term is not penalized.

Including these additional terms encourages smoother control-input and
flow-velocity sequences and avoids unnecessarily large changes between
consecutive optimization steps.


\subsection{Numerical Solver Change}

The numerical solver used for the nonlinear optimization problem was also
changed during the implementation.

The earlier stochastic bilinear models were solved using Sequential Least
Squares Programming (SLSQP). As the district heating model was extended,
the optimization problem became more computationally demanding. In
particular, the larger state dimension, propagation of the stochastic
covariance, chance constraints, bilinear dependence on \(V\), and the
additional penalties on the changes in \(U\) and \(V\) increased the
computation time.

For this reason, the implementation was changed from SLSQP to IPOPT. The
optimization problem is constructed using CasADi.

\section{Results}
\label{sec:district_heating_results}

The developed controller was evaluated on the single-pipe district heating
example. The pipe was divided into \(M=6\) finite volumes over a total length
of \(6\,\mathrm{km}\), and the sampling time was selected as
\(\Delta t=0.5\,\mathrm{h}\). The heat-loss coefficient was

\begin{equation}
    \beta = 7.66\times10^{-3}\,\mathrm{h}^{-1}
\end{equation}

Two temperatures were considered as measured outputs: the temperature at the
pipe outlet and the temperature at the third finite volume. The initial
temperature of the pipe was \(65^\circ\mathrm{C}\), while the initial ambient
temperature was \(0^\circ\mathrm{C}\).

The physical output constraints were selected as

\begin{equation}
    65^\circ\mathrm{C}
    \leq
    T_{\mathrm{out}}
    \leq
    90^\circ\mathrm{C}
\end{equation}

for both measured locations. The supply-temperature and flow-velocity bounds
were

\begin{equation}
    50^\circ\mathrm{C}
    \leq
    T_s
    \leq
    110^\circ\mathrm{C}
\end{equation}

and

\begin{equation}
    0.05
    \leq
    v
    \leq
    7.0
    \quad \mathrm{km/h}
\end{equation}

The reference temperatures were changed during the simulation according to

\begin{equation}
    \bar{T}_1
    =
    \begin{bmatrix}
        70 & 70
    \end{bmatrix}^{T},
    \qquad
    \bar{T}_2
    =
    \begin{bmatrix}
        70 & 87
    \end{bmatrix}^{T},
    \qquad
    \bar{T}_3
    =
    \begin{bmatrix}
        87 & 70
    \end{bmatrix}^{T}
\end{equation}

Each reference level was held for \(48\) sampling instants, corresponding to
\(24\,\mathrm{h}\). The ambient temperature was initially
\(0^\circ\mathrm{C}\) and was changed to \(10^\circ\mathrm{C}\) at
\(t=36\,\mathrm{h}\).

The weighting matrices used in the optimization were

\begin{equation}
    Q=10^{11}I,
    \qquad
    R_U=0,
    \qquad
    R_V=0
\end{equation}

together with the move-suppression weights

\begin{equation}
    R_{\Delta U}=10,
    \qquad
    R_{\Delta V}=20
\end{equation}

The stochastic behavior was evaluated using \(200\) Monte Carlo
realizations.


\subsection{Open-Loop}
\label{subsec:district_open_loop_results}

In the open-loop simulation, the optimization problem was solved once over
the complete simulation horizon. The complete future ambient-temperature
profile was provided to the controller before the optimization. Therefore,
the calculated sequences \(U^\ast\) and \(V^\ast\) already account for the
scheduled change in ambient temperature.

The same optimized input and velocity sequences were subsequently applied to
all Monte Carlo realizations. 




% Example figure:
%
% \begin{figure}[H]
%     \centering
%     \includegraphics[width=\textwidth]
%     {open_loop_physical_summary.png}
%     \caption{Open-loop Monte Carlo results for the single-pipe district
%     heating system.}
%     \label{fig:district_open_loop}
% \end{figure}

% \begin{figure}[H]
%     \centering
%     \includegraphics[width=\textwidth]
%     {open_loop_z_profile.png}
%     \caption{Open-loop temperature propagation along the pipe.}
%     \label{fig:district_open_loop_profile}
% \end{figure}


\subsection{Closed-Loop}
\label{subsec:district_closed_loop_results}

For the closed-loop simulation, a prediction horizon of

\begin{equation}
    N=12
\end{equation}

was used, corresponding to \(6\,\mathrm{h}\). In contrast to the open-loop
case, the controller does not receive the future ambient-temperature
trajectory. At each sampling instant, only the currently available ambient
temperature is used and is assumed to remain constant over the prediction
horizon.

A separate MPC optimization is performed for every Monte Carlo realization.
Consequently, each trajectory obtains its own optimal \(U^\ast\) and
\(V^\ast\) sequences according to its current state estimate. Only the first
elements of these sequences are applied, after which a new measurement is
received and the optimization is repeated.

When a new ambient-temperature value becomes available, the Kalman-filter
state prediction is corrected using the ambient-temperature re-basing term.
The corrected state estimate is then used by the subsequent MPC problem.
This allows the controller to respond to the actual ambient-temperature
change without assuming its future value beforehand.



% Example figure:
%
% \begin{figure}[H]
%     \centering
%     \includegraphics[width=\textwidth]
%     {closed_loop_physical_summary.png}
%     \caption{Closed-loop Monte Carlo results for the single-pipe district
%     heating system.}
%     \label{fig:district_closed_loop}
% \end{figure}

% \begin{figure}[H]
%     \centering
%     \includegraphics[width=\textwidth]
%     {closed_loop_z_profile.png}
%     \caption{Closed-loop temperature propagation along the pipe.}
%     \label{fig:district_closed_loop_profile}
% \end{figure}